% Minimal preamble for Open Logic Project sources (pinned commit 1e960beff9ed7835bf3e3f1335e21af3439cd107).
%
% OLP's own master preamble cannot run on the project host: TeX Live 2025 there lacks
% prooftree.sty, ntheorem.sty, extarrows.sty, mathdots.sty, isomath.sty and subfiles.sty.
% Following the project's own rule ("if the author preamble needs packages the toolchain lacks,
% do not invent macros -- use a minimal working preamble"), this file declares, verbatim in
% spirit and numbering behaviour, exactly the theorem-like environments OLP defines in
% sty/open-logic-book-envs.sty (thm/lem/prop/cor/defn/ex/prob/rem/note/case/axiom/conv),
% plus the shorthand accents OLP uses (`!!a`, `!!b`, `!!{...}`).
%
% Author text is untouched; only the packaging is ours.
\documentclass[10pt,a4paper]{article}

\usepackage{amsmath}
\usepackage{amssymb}
\usepackage{amsthm}
\usepackage{graphicx}
\usepackage{hyperref}
\usepackage{enumitem}
\usepackage[framemethod=default]{mdframed}

% ---- OLP shorthand accents (from sty/open-logic.sty) ----
% NB: \! is taken by LaTeX (negative thin space), so we only define the !!a / !!b families
% and the \olup helper used by them.
\newcommand*{\olup}[1]{\MakeUppercase{#1}}
\def\!!!a#1{\olup{#1}}
\def\!!a#1{\olup{#1}}
\def\!!b#1{\olup{#1}}

% ---- theorem-like environments, same family numbering as OLP ----
\theoremstyle{plain}
\newtheorem{thm}{Theorem}[section]
\newtheorem{lem}[thm]{Lemma}
\newtheorem{prop}[thm]{Proposition}
\newtheorem{cor}[thm]{Corollary}
\newtheorem{axiom}{Axiom}

\theoremstyle{definition}
\newtheorem{defn}[thm]{Definition}
\newtheorem{ex}[thm]{Example}
\newtheorem{prob}{Problem}[section]
\newtheorem{probd}{Problem}
\newtheorem{conv}[thm]{Convention}

\theoremstyle{remark}
\newtheorem{rem}[thm]{Remark}
\newtheorem{note}[thm]{Note}
\newtheorem{case}[thm]{Case}

% ---- environments OLP defines as wrappers ----
\newenvironment{probtag}[1]{\begin{prob}}{\end{prob}}
\newenvironment{explain}{}{}
\newenvironment{explain*}{}{}
\newenvironment{editorial}{}{}
\newenvironment{digress}{}{}
\newenvironment{digress*}{}{}
\newenvironment{history}{}{}
\newenvironment{reference}{}{}
\newcommand*{\olfileid}[3]{}
\newcommand*{\olsection}[1]{\section*{#1}}
\newcommand*{\olchapter}[3]{}
\newcommand*{\olpart}[2]{}
\newcommand*{\olimport}[1]{}
\newcommand*{\olinclude}[1]{}
\newcommand*{\olpath}[1]{}
\newcommand*{\ollangid}[1]{}
\newcommand*{\problemtag}[1]{}
\newcommand*{\problemauthor}[1]{}
\newcommand*{\ollabel}[1]{\label{#1}}
% \olref takes 0-3 OPTIONAL arguments plus one mandatory label:
%     \olref{lab}  |  \olref[from]{lab}  |  \olref[from][to]{lab}  |  \olref[a][b][c]{lab}
% A single-argument definition silently swallows "[sth]" and emits a broken reference
% (measured on the host: the text came out as "??iso]x"). The chain below consumes the
% brackets and cross-references only the mandatory argument.
\makeatletter
\newcommand{\olref}{\@ifnextchar[{\olref@i}{\olref@final}}
\def\olref@i[#1]{\@ifnextchar[{\olref@ii}{\olref@final}}
\def\olref@ii[#1]{\@ifnextchar[{\olref@iii}{\olref@final}}
\def\olref@iii[#1]{\olref@final}
\def\olref@final#1{\ref{#1}}
\makeatother
\newcommand*{\olnewpage}{\clearpage}

% ===================== OLP compatibility layer (1) =====================

\makeatletter

% ---------------- proof trees (OLP's prooftree is a bussproofs derivative) ----------------
\providecommand{\AxiomC}[1]{}
\providecommand{\UnaryInfC}[1]{}
\providecommand{\BinaryInfC}[1]{}
\providecommand{\TrinaryInfC}[1]{}
\providecommand{\QuaternaryInfC}[1]{}
\providecommand{\QuinaryInfC}[1]{}
\providecommand{\RightLabel}[1]{}
\providecommand{\LeftLabel}[1]{}
\providecommand{\doubleLine}{}
\providecommand{\noLine}{}
\providecommand{\DisplayProof}{}
\providecommand{\fCenter}{}
\newenvironment{prooftree}{}{}
\newenvironment{ltprooftree}{}{}
\newenvironment{oltableau}{}{}

% ---------------- logical notation ----------------
\providecommand{\lif}{\mathbin{\rightarrow}}
\providecommand{\lforall}{\mathord{\forall}}
\providecommand{\lexists}{\mathord{\exists}}
\providecommand{\ltrue}{\top}
\providecommand{\lfalse}{\bot}
\providecommand{\Ltrue}{\top}
\providecommand{\Lfalse}{\bot}
\providecommand{\Sequent}{\mathrel{\Rightarrow}}
\providecommand{\Proves}{\mathrel{\vdash}}
\providecommand{\Entails}{\mathrel{\models}}
\providecommand{\Subst}[3]{#1[#2/#3]}
\providecommand{\sFmla}[2]{#1(#2)}
\providecommand{\pFmla}[2]{#1(#2)}
\providecommand{\formula}{\relax}
\providecommand{\True}{\top}
\providecommand{\False}{\bot}
\providecommand{\Args}[1]{#1}
\providecommand{\Ts}{\ensuremath{\top}}
\providecommand{\Fi}{\ensuremath{\bot}}
\providecommand{\lfalsum}{\bot}

% ---------------- objects, structures, metatheory ----------------
\providecommand{\Obj}{}
\providecommand{\Struct}{\ensuremath{\mathfrak{A}}}
\providecommand{\Domain}{\ensuremath{|\mathfrak{A}|}}
\providecommand{\Assign}{\ensuremath{s}}
\providecommand{\Sat}[2]{\ensuremath{#1 \models #2}}
\providecommand{\mSat}[3]{\ensuremath{#1 \models #2[#3]}}
\providecommand{\mModel}[2]{\ensuremath{#1 \models #2}}
\providecommand{\Lang}{\ensuremath{\mathcal{L}}}
\providecommand{\Th}{\ensuremath{\mathrm{Th}}}
\providecommand{\Log}{\ensuremath{\mathrm{L}}}
\providecommand{\tuple}[1]{\langle #1 \rangle}
\providecommand{\num}[1]{\underline{#1}}
\providecommand{\cfind}[1]{\varphi_{#1}}
\providecommand{\Nat}{\ensuremath{\mathbb{N}}}
\providecommand{\Mult}{\ensuremath{\mathrm{Mult}}}
\providecommand{\Add}{\ensuremath{\mathrm{Add}}}
\providecommand{\Exp}{\ensuremath{\mathrm{Exp}}}
\providecommand{\Zero}{\ensuremath{0}}
\providecommand{\Succ}{\ensuremath{\mathrm{S}}}
\providecommand{\lambd}{\ensuremath{\lambda}}
\providecommand{\bredpar}{\mathrel{\rightarrow_\beta}}
\providecommand{\bcd}[1]{\ensuremath{\ulcorner #1 \urcorner}}
\providecommand{\Set}[1]{\ensuremath{\{#1\}}}
\providecommand{\indcase}[1]{\par\noindent\emph{Case #1}.\quad}
\providecommand{\indfrm}[1]{\par\noindent\emph{Form #1}.\quad}
\providecommand{\iftag}[1]{#1}
\providecommand{\tagitem}[1]{\item[#1]}
\providecommand{\Elim}[1]{}
\providecommand{\Intro}[1]{}
\providecommand{\LeftR}[1]{}
\providecommand{\RightR}[1]{}
\providecommand{\TAss}[1]{}
\providecommand{\TRule}[1]{}

\makeatother
\providecommand{\olpath}{.}

% ===================== OLP compatibility layer (2) =====================
% --- second compatibility block: notation and shorthands used by the frozen OLP units ---
\providecommand{\CutCS}{\mathrm{Cut}}
\providecommand{\LeftSubproofLabel}{\relax}
\providecommand{\RightSubproofLabel}{\relax}
\providecommand{\maeh}{\relax}
\providecommand{\cheight}{\relax}
\providecommand{\ande}{\mathrm{and}_E}
\providecommand{\andi}{\mathrm{and}_I}
\providecommand{\ore}{\mathrm{or}_E}
\providecommand{\ori}{\mathrm{or}_I}
\providecommand{\typeof}{\mathord{:}}
\providecommand{\pValue}[1]{#1}
\providecommand{\Sent}{\mathrm{Sent}}
\providecommand{\cardplus}{\mathbin{+}}
\providecommand{\cardtimes}{\mathbin{\cdot}}
\providecommand{\cardexpo}{\mathbin{\hat{\ }}}
\providecommand{\cardnless}{\not<}
\providecommand{\cardle}{\leq}
\providecommand{\funfromto}[3]{#1\colon #2 \to #3}
\providecommand{\funimage}[1]{\operatorname{im}(#1)}
\providecommand{\fdefined}{\mathord{\downarrow}}
\providecommand{\canonord}{\operatorname{ord}}
\providecommand{\ordtimes}{\mathbin{\cdot}}
\providecommand{\ordeq}[2]{#1 = #2}
\providecommand{\ordneq}[2]{#1 \neq #2}
\providecommand{\ordtype}[1]{\operatorname{ord}(#1)}
\providecommand{\ordsucc}{\mathord{+1}}
\providecommand{\ordunion}{\cup}
\providecommand{\onesphere}{S^1}
\providecommand{\rotationsgroup}{\mathrm{SO}(2)}
\providecommand{\equivrep}[1]{[#1]}
\providecommand{\rlexless}{\prec}
\providecommand{\ZFC}{\mathrm{ZFC}}
\providecommand{\ZF}{\mathrm{ZF}}
\providecommand{\Rat}{\mathbb{Q}}
\providecommand{\Real}{\mathbb{R}}
\providecommand{\PosInt}{\mathbb{Z}^{+}}
\providecommand{\Id}{\mathrm{id}}
\providecommand{\Bin}{\mathrm{Bin}}
\providecommand{\pto}{\rightharpoonup}
\providecommand{\dom}{\operatorname{dom}}
\providecommand{\ran}{\operatorname{ran}}
\providecommand{\len}{\operatorname{len}}
\providecommand{\concat}{\mathbin{{}^\frown}}
\providecommand{\card}[1]{\lvert #1 \rvert}
\providecommand{\Pow}{\mathcal{P}}
\providecommand{\Setabs}[2]{\{\, #1 \mid #2 \,\}}
\providecommand{\comp}{\circ}
\providecommand{\Char}{\chi}
\providecommand{\Prf}{\operatorname{Prf}}
\providecommand{\Prov}{\operatorname{Prov}}
\providecommand{\Proj}{\operatorname{proj}}
\providecommand{\Gn}{\operatorname{gn}}
\providecommand{\gn}{\operatorname{gn}}
\providecommand{\Value}[1]{\operatorname{val}(#1)}
\providecommand{\Deduce}{\mathrm{D}}
\providecommand{\Dual}{\mathrm{d}}
\providecommand{\Ax}{\mathrm{Ax}}
\providecommand{\pSat}[3]{#1 \models #2\,[#3]}
\providecommand{\pAssign}{\sigma}
\providecommand{\bforall}{\forall}
\providecommand{\bexists}{\exists}
\providecommand{\liff}{\mathbin{\leftrightarrow}}
\providecommand{\FV}{\operatorname{FV}}
\providecommand{\UnaryInf}{\ensuremath{\mid}}
\providecommand{\BinaryInf}{\ensuremath{\mid}}
\providecommand{\depth}{\operatorname{depth}}
\providecommand{\oliflabeldef}[3]{#1}
\providecommand{\TMblank}{\sqcup}
\providecommand{\TMstroke}{\mid}
\providecommand{\TMstay}{\mathrm{S}}
\providecommand{\TMright}{\mathrm{R}}
\providecommand{\TMleft}{\mathrm{L}}
\providecommand{\TMendtape}{\triangleright}
\providecommand{\umin}{\mu}
\providecommand{\Sum}{\sum}
\providecommand{\citep}[1]{\cite{#1}}
\providecommand{\citet}[1]{\cite{#1}}
\providecommand{\citeauthor}[1]{\cite{#1}}
\providecommand{\citeyear}[1]{\cite{#1}}

% ===================== OLP compatibility layer (3) =====================
% --- third compatibility block: remaining notation seen in the frozen OLP units ---
\providecommand{\iso}{\cong}
\providecommand{\nszero}{\mathrm{0}}
\providecommand{\nssucc}{\mathrm{S}}
\providecommand{\nsplus}{\mathbin{+}}
\providecommand{\nssless}{<}
\providecommand{\nsless}{<}
\providecommand{\PL}{\mathrm{PL}}
\providecommand{\Nec}{\mathord{\Box}}
\providecommand{\RK}{\mathrm{RK}}
\providecommand{\ST}{\mathrm{ST}}
\providecommand{\SSubst}[3]{#1[#2/#3]}
\providecommand{\subst}[3]{#1[#2/#3]}
\providecommand{\tagfalse}{\mathrm{F}}
\providecommand{\cutr}{\mathrm{Cut}}
\providecommand{\cardneq}{\neq}
\providecommand{\issucc}[1]{\operatorname{succ}(#1)}
\providecommand{\tup}[1]{\langle #1 \rangle}
\providecommand{\Var}{\mathrm{Var}}
\providecommand{\Bew}{\operatorname{Bew}}
\providecommand{\Con}{\operatorname{Con}}
\providecommand{\Cons}{\operatorname{Cons}}
\providecommand{\pDeduce}{\mathrm{D}}
\providecommand{\psat}{\models}
\providecommand{\den}[1]{\llbracket #1 \rrbracket}
\providecommand{\Sem}[2]{\llbracket #1 \rrbracket_{#2}}
\providecommand{\DomainOf}[1]{\lvert #1 \rvert}
\providecommand{\Form}{\mathrm{Form}}
\providecommand{\Terms}{\mathrm{Term}}
\providecommand{\Atom}{\mathrm{Atom}}
\providecommand{\FVars}[1]{\operatorname{FV}(#1)}
\providecommand{\fv}[1]{\operatorname{fv}(#1)}
\providecommand{\tuplep}[1]{\langle #1 \rangle}
\providecommand{\restrict}[2]{#1 \upharpoonright #2}
\providecommand{\sto}{\rightarrow}
\providecommand{\FmlaSet}[1]{\mathrm{Fmla}_{#1}}
\providecommand{\Thm}{\mathrm{Thm}}
\providecommand{\Eval}{\operatorname{val}}
\providecommand{\blank}{\sqcup}
\providecommand{\blankTape}{\sqcup}
\providecommand{\Consis}{\operatorname{Con}}
\providecommand{\drs}{\mathrm{d}}
\providecommand{\VRule}{\mid}
\providecommand{\Bet}{\mathrm{B}}
\providecommand{\Seg}[1]{\mathrm{Seg}_{#1}}
\providecommand{\Enums}{\mathrm{En}}
\providecommand{\Z}{\mathbb{Z}}
\providecommand{\N}{\mathbb{N}}
\providecommand{\abs}[1]{\lvert #1 \rvert}
\providecommand{\floor}[1]{\lfloor #1 \rfloor}
\providecommand{\ceil}[1]{\lceil #1 \rceil}

% ===================== OLP compatibility layer (4) =====================
% --- fourth compatibility block: harvested from the per-candidate compile logs ---
\providecommand{\beredpar}{\mathrel{\rightarrow_\beta}}
\makeatletter
\def\bered{\@ifnextchar[{\bered@i}{\mathrel{\rightarrow_\beta}}}
\def\bered@i[#1]{\mathrel{\rightarrow_\beta}}
\def\bred{\@ifnextchar[{\bred@i}{\mathrel{\rightarrow_\beta}}}
\def\bred@i[#1]{\mathrel{\rightarrow_\beta}}
\makeatother
\providecommand{\beredone}{\mathrel{=_\beta}}
\providecommand{\bredone}{\mathrel{=_\beta}}
\providecommand{\xredone}{\mathrel{=_\beta}}
\providecommand{\aconvone}{\mathrel{=_\alpha}}
\providecommand{\cardless}{<}
\providecommand{\tagrefs}[1]{#1}
\providecommand{\ordplus}{\mathbin{+}}
\providecommand{\funrestrictionto}[2]{#1 \upharpoonright #2}
\providecommand{\fundefined}{\mathord{\uparrow}}
\providecommand{\Trm}{\mathrm{Trm}}
\providecommand{\OProv}{\operatorname{Prov}}
\providecommand{\OPrf}{\operatorname{Prf}}
\providecommand{\aeq}{\mathrel{=}}
\providecommand{\mClass}[1]{[#1]}
\providecommand{\equivclass}[1]{[#1]}
\providecommand{\pheight}{\operatorname{ht}}
\providecommand{\phy}{\operatorname{ht}}
\providecommand{\Cut}{\mathrm{Cut}}
\providecommand{\QuantRank}{\operatorname{qr}}
\providecommand{\elemequiv}{\mathrel{\equiv}}
\providecommand{\PIso}{\mathrm{PIso}}
\providecommand{\Theory}{\mathrm{Th}}
\providecommand{\supstrict}{\subsetneq}
\providecommand{\trcl}{\operatorname{trcl}}
\providecommand{\cref}[1]{\ref{#1}}
\providecommand{\Discharge}{\mathrm{Dis}}
\providecommand{\Weakening}{\mathrm{Wk}}
\providecommand{\ZFminus}{\mathrm{ZF}^{-}}
\providecommand{\RK}{\mathrm{RK}}
\makeatletter
\renewcommand{\Dual}{\@ifnextchar[{\Dual@i}{\ifmmode\mathrm{d}\else$\mathrm{d}$\fi}}
\def\Dual@i[#1]{\ifmmode\mathrm{d}[#1]\else$\mathrm{d}[#1]$\fi}
\makeatother
\providecommand{\RightR}{\mathrm{R}}
\providecommand{\LeftR}{\mathrm{L}}
\providecommand{\Intro}{\mathrm{I}}
\providecommand{\Axiom}{\mathrm{Ax}}
\providecommand{\formula}[1]{#1}
\providecommand{\stared}{\ast}
\providecommand{\oset}[2]{\{\, #1 \mid #2 \,\}}
\providecommand{\setof}[2]{\{\, #1 \mid #2 \,\}}
\providecommand{\pow}[1]{\mathcal{P}(#1)}
\providecommand{\ordpair}[2]{\langle #1, #2 \rangle}
\providecommand{\ordtriple}[3]{\langle #1, #2, #3 \rangle}
\providecommand{\union}{\cup}
\providecommand{\inter}{\cap}
\providecommand{\seq}[1]{\langle #1 \rangle}
\providecommand{\SuchThat}{\mid}
\providecommand{\where}{\mid}
\providecommand{\TVar}{\mathrm{TVar}}
\providecommand{\Bool}{\mathrm{Bool}}
\providecommand{\Val}{\operatorname{Val}}
\providecommand{\FV}{\operatorname{FV}}
\providecommand{\TL}[1]{\mathrm{TL}}
\providecommand{\xstared}{\ast}

% ===================== OLP compatibility layer (5) =====================
% --- fifth compatibility block: rescues the quarantined OLP candidates ---
\providecommand{\ident}{\operatorname{id}}
\providecommand{\setrank}{\operatorname{rank}}
\providecommand{\cardfont}{\mathfrak{c}}
\providecommand{\Complement}{\overline}
\providecommand{\QR}{\mathrm{QR}}
\providecommand{\becd}{\ensuremath{\ulcorner}}
\providecommand{\disjointsum}{\mathbin{\oplus}}
\providecommand{\cardeq}[2]{#1 = #2}
\providecommand{\xred}{\mathrel{\rightarrow}}
\providecommand{\cardless}{<}
\providecommand{\cardplus}{\mathbin{+}}
\providecommand{\cardtimes}{\mathbin{\cdot}}
\providecommand{\cardexpo}{\mathbin{\hat{\ }}}
\providecommand{\cardneq}{\neq}
\providecommand{\cardle}{\leq}
\providecommand{\cardnless}{\not<}
% --- sixth compatibility block: names reported by real compile logs ---
\providecommand{\ORefut}{\operatorname{Refut}}
\providecommand{\ORProv}{\operatorname{Prov}}
\providecommand{\cutrank}{\operatorname{cr}}
\providecommand{\Refut}{\operatorname{Refut}}
\providecommand{\DischargeRule}{\mathrm{Dis}}
\providecommand{\Contraction}{\mathrm{C}}
\providecommand{\bottomAlignProof}{\relax}
\providecommand{\FalseCl}{\bot}
\providecommand{\red}{\mathrel{\rightarrow}}
\providecommand{\varAssign}{\sigma}
\providecommand{\defiff}{\mathrel{\leftrightarrow}}
\providecommand{\FalseInt}{\bot}
\providecommand{\fn}{\operatorname{fn}}
\providecommand{\stared}{\ast}
\providecommand{\dotted}{\cdot}
\providecommand{\Gamma}{\Gamma}
% --- sixth compatibility block: names reported by real compile logs ---
\providecommand{\ORefut}{\operatorname{Refut}}
\providecommand{\ORProv}{\operatorname{Prov}}
\providecommand{\cutrank}{\operatorname{cr}}
\providecommand{\Refut}{\operatorname{Refut}}
\providecommand{\DischargeRule}{\mathrm{Dis}}
\providecommand{\Contraction}{\mathrm{C}}
\providecommand{\bottomAlignProof}{\relax}
\providecommand{\FalseCl}{\bot}
\providecommand{\red}{\mathrel{\rightarrow}}
\providecommand{\varAssign}{\sigma}
\providecommand{\defiff}{\mathrel{\leftrightarrow}}
\providecommand{\FalseInt}{\bot}
\providecommand{\fn}{\operatorname{fn}}
\providecommand{\stared}{\ast}
\providecommand{\dotted}{\cdot}
\providecommand{\Gamma}{\Gamma}
\begin{document}

% ==== OLP file: content/set-theory/cardinals/cardinals.tex ====
% Part: set-theory
% Chapter: cardinals



\olchapter{sth}{cardinals}{Cardinals}

\olimport{cp}
\olimport{cardsasords}
\olimport{milestone}
\olimport{classing}
\olimport{hp}

\OLEndChapterHook

% ==== OLP file: content/set-theory/cardinals/cardsasords.tex ====


\olfileid{sth}{cardinals}{cardsasords}
\olsection{Cardinals as Ordinals}

In fact, our theory of cardinals will just make (shameless) use of our
theory of ordinals. That is: we will just define cardinals as certain
specific ordinals. In particular, we will offer the following:

\begin{defn}\ollabel{defcardinalasordinal}
If $A$ can be well-ordered, then $\card{A}$ is the least ordinal
$\gamma$ such that $\cardeq{A}{\gamma}$. For any ordinal $\gamma$, we
say that $\gamma$ is a \emph{cardinal} iff $\gamma = \card{\gamma}$.
\end{defn}

We just used the phrase ``$A$ can be well-ordered''. As is almost
always the case in mathematics, the modal locution here is just a
hand-waving gloss on an existential claim: to say ``$A$ can be
well-ordered'' is just to say ``there is a relation which
well-orders~$A$''. 

But there is a snag with \olref{defcardinalasordinal}. We would like
it to be the case that \emph{every} set has a size, i.e., that
$\card{A}$ exists for every~$A$. The definition we just gave, though,
begins with a conditional: ``\emph{If} $A$ can be
well-ordered\ldots''. If there is some set $A$ which cannot be
well-ordered, then our definition will simply fail to define an object~$\card{A}$.

So, to use \olref{defcardinalasordinal}, we need a guarantee that
every set can be well-ordered. Sadly, though, this guarantee is
unavailable in~$\ZF$. So, if we want to use
\olref{defcardinalasordinal}, there is no alternative but to add a new
axiom, such as:
\begin{axiom}[Well-Ordering]
Every set can be well-ordered.
\end{axiom}
We will discuss whether the Well-Ordering Axiom is acceptable in
\olref[choice][]{chap}. From now on, though, we will simply help
ourselves to it. And, using it, it is quite straightforward to prove
that cardinals (as defined in \olref{defcardinalasordinal}) exist and
behave nicely:

\begin{lem}\ollabel{lem:CardinalsExist}
For every set $A$:
\begin{enumerate}
	\item\ollabel{cardaexists} $\card{A}$ exists and is unique;
	\item\ollabel{cardaapprox}  $\cardeq{\card{A}}{A}$;
	\item\ollabel{cardaidem}  $\card{A}$ is a cardinal, i.e.,
	$\card{A} = \card{\card{A}}$;
\end{enumerate}
\end{lem}

\begin{proof}
Fix $A$. By Well-Ordering, there is a well-ordering $\tuple{A, R}$. By
\olref[ordinals][ordtype]{thmOrdinalRepresentation}, $\tuple{A,
R}$ is isomorphic to a unique ordinal, $\beta$. So
$\cardeq{A}{\beta}$. By Transfinite Induction, there is a uniquely
least ordinal, $\gamma$, such that $\cardeq{A}{\gamma}$. So $\card{A}
= \gamma$, establishing \olref{cardaexists} and \olref{cardaapprox}.
To establish \olref{cardaidem}, note that if $\delta \in \gamma$ then
$\cardless{\delta}{A}$, by our choice of $\gamma$, so that also
$\cardless{\delta}{\gamma}$ since equinumerosity is an equivalence
relation (\olref[sfr][siz][equ]{equinumerosityisequi}). So $\gamma =
\card{\gamma}$. 
%So, for reductio, suppose that there is some ordinal $\delta \in \gamma$ such that $\cardeq{\gamma}{\delta}$. Then, $\cardeq{\cardeq{A}{\gamma}}{\delta}$ so that $\cardeq{A}{\delta}$ by \olref[sfr][set][equ]{equinumerosityisequi}, which contradicts the choice of $\gamma$ as the least ordinal such that $\cardeq{A}{\gamma}$. 
\end{proof}

The next result guarantees Cantor's Principle, and more besides.
(Note that cardinals inherit their ordering from the ordinals, i.e.,
$\cardfont{a} < \cardfont{b}$ iff $\cardfont{a} \in \cardfont{b}$. In
formulating this, we will use Fraktur letters for objects we know to be
cardinals. This is fairly standard. A common alternative is to use
Greek letters, since cardinals are ordinals, but to choose them from
the middle of the alphabet, e.g.: $\kappa, \lambda$.):
\begin{lem}\ollabel{lem:CardinalsBehaveRight}
For any sets $A$ and $B$:
\begin{align*}
	\cardeq{A}{B} &\text{ iff } \card{A} = \card{B}\\
	\cardle{A}{B} &\text{ iff } \card{A} \leq \card{B}\\
	\cardless{A}{B}&\text{ iff } \card{A} < \card{B}
\end{align*}
\end{lem}

\begin{proof}
We will prove the left-to-right direction of the second claim (the
other cases are similar, and left as an exercise). So, consider the
following diagram:
\begin{center}
	\begin{tikzpicture}
	\node (nodea) {$A$};
	\node[right = 6em of nodea] (nodeb) {$B$};
	\node[below = 2em of nodea] (nodecarda) {$\card{A}$};
	\node[below = 2em of nodeb] (nodecardb) {$\card{B}$};
	\draw[->] (nodea)--(nodeb);
	\draw[<->] (nodea)--(nodecarda);
	\draw[<->] (nodeb)--(nodecardb);
	\draw[->, dashed] (nodecarda)--(nodecardb);
\end{tikzpicture}
\end{center}
The double-headed arrows indicate !!{bijection}s, whose existence is
guaranteed by \olref{lem:CardinalsExist}. In assuming that
$\cardle{A}{B}$, there is !!a{injection}  $A\to B$. Now,
chasing the arrows around from $\card{A}$ to $A$ to $B$ to $\card{B}$,
we obtain !!a{injection} $\card{A} \to \card{B}$ (the dashed arrow).
\end{proof}\noindent We can also use \olref{lem:CardinalsBehaveRight}
to re-prove Schr\"{o}der--Bernstein. This is the claim that if
$\cardle{A}{B}$ and $\cardle{B}{A}$ then $\cardeq{A}{B}$. We stated
this as \olref[sfr][siz][sb]{thm:schroder-bernstein}, but first proved
it---with some effort---in \olref[sfr][infinite][card-sb]{sec}.
Now consider:

\begin{proof}[Re-proof of Schr\"oder-Bernstein]
If $\cardle{A}{B}$ and $\cardle{B}{A}$, then $\card{A} \leq \card{B}$
and $\card{B} \leq \card{A}$ by \olref{lem:CardinalsBehaveRight}. So
$\card{A} = \card{B}$ and $\cardeq{A}{B}$ by Trichotomy and
\olref{lem:CardinalsBehaveRight}.
\end{proof}
\noindent
Whilst this is a very simple proof, it implicitly relies on both
Replacement (to secure
\olref[ordinals][ordtype]{thmOrdinalRepresentation}) and on
Well-Ordering (to guarantee \olref{lem:CardinalsBehaveRight}). By
contrast, the proof of \olref[sfr][infinite][card-sb]{sec} was much
more self-standing (indeed, it can be carried out in~$\Zminus$).

% ==== OLP file: content/set-theory/cardinals/classing.tex ====


\olfileid{sth}{cardinals}{classing}
\olsection{Finite, \usetoken{S}{enumerable}, \usetoken{S}{nonenumerable}}

Now that we have been introduced to cardinals, it is worth spending a
little time talking about different varieties of cardinals;
specifically, finite, !!{enumerable}, and !!{nonenumerable} cardinals.

Our first two results entail that the finite cardinals will be exactly
the finite ordinals, which we defined as our \emph{natural numbers}
back in \olref[z][infinity-again]{defnomega}: 

\begin{prop}\ollabel{finitecardisoequal}
Let $n, m \in \omega$. Then $n = m$ iff $\cardeq{n}{m}$.
\end{prop}

\begin{proof}
\emph{Left-to-right} is trivial. To prove \emph{right-to-left},
suppose $\cardeq{n}{m}$ although $n \neq m$. By Trichotomy, either $n
\in m$ or $m \in n$; suppose $n \in m$ without loss of generality.
Then $n \subsetneq m$ and there is !!a{bijection} $f \colon m \to n$,
so that $m$ is Dedekind infinite, contradicting
\olref[z][infinity-again]{naturalnumbersarentinfinite}.
\end{proof}

\begin{cor}\ollabel{naturalsarecardinals}
If $n \in \omega$, then $n$ is a cardinal. 
\end{cor}

\begin{proof}
Immediate.
\end{proof}
\noindent
It also follows that several reasonable notions of what it might mean
to describe a cardinal as ``finite'' or ``infinite'' coincide:
\begin{thm}\ollabel{generalinfinitycharacter}For any set $A$, the following are equivalent:
\begin{enumerate}
	\item\ollabel{card:notinomega} $\card{A} \notin \omega$, i.e.,
	$A$ is not a natural number;
	\item\ollabel{card:omegaplus} $\omega \leq \card{A}$;
	\item\ollabel{card:infinite} $A$ is Dedekind infinite.
\end{enumerate}
\end{thm}

\begin{proof}
From \olref[ord-arithmetic][using-addition]{ordinfinitycharacter},
\olref[cardinals][cardsasords]{lem:CardinalsBehaveRight}, and
\olref{naturalsarecardinals}. 
\end{proof}

This licenses the following \emph{definition} of some notions which we
used rather informally in \olref[sfr][][]{part}:

\begin{defn}\ollabel{defnfinite}
We say that $A$ is \emph{finite} iff $\card{A}$ is a natural number,
i.e., $\card{A} \in \omega$. Otherwise, we say that $A$ is
\emph{infinite}.
\end{defn}
\noindent 
But note that this definition is presented against the background of
$\ZFC$. After all, we needed Well-Ordering to guarantee that every set
has a cardinality. And indeed, without Well-Ordering, there can be a
set which is neither finite nor Dedekind infinite. We will return to
this sort of issue in \olref[choice][]{chap}. For now, we continue to
rely upon Well-Ordering.

Let us now turn from the finite cardinals to the infinite cardinals.
Here are two elementary points:

\begin{cor}\ollabel{omegaisacardinal}
$\omega$ is the least infinite cardinal. 
\end{cor}

\begin{proof}
$\omega$ is a cardinal, since $\omega$ is Dedekind infinite and if
$\cardeq{\omega}{n}$ for any $n \in \omega$ then $n$ would be Dedekind
infinite, contradicting
\olref[z][infinity-again]{naturalnumbersarentinfinite}. Now
$\omega$ is the least infinite cardinal by definition. 
\end{proof}

\begin{cor}
Every infinite cardinal is a limit ordinal.
\end{cor}

\begin{proof}
Let $\alpha$ be an infinite successor ordinal, so $\alpha = \beta
\ordplus 1$ for some $\beta$. By \olref{finitecardisoequal}, $\beta$
is also infinite, so $\cardeq{\beta}{\beta \ordplus 1}$ by
\olref[ord-arithmetic][using-addition]{ordinfinitycharacter}. Now
$\card{\beta} = \card{\beta\ordplus 1} = \card{\alpha}$ by
\olref[cardinals][cardsasords]{lem:CardinalsBehaveRight}, so
that $\alpha \neq \card{\alpha}$.
\end{proof}

Now, as early as \olref[sfr][siz][enm-alt]{defn:enumerable}, we flagged we
can distinguish between !!{enumerable} and !!{nonenumerable} infinite
sets. That definition naturally leads to the following:

\begin{prop}
$A$ is !!{enumerable} iff $\card{A} \leq \omega$, and $A$ is
!!{nonenumerable} iff $\omega < \card{A}$.
\end{prop}

\begin{proof}
By Trichotomy, the two claims are equivalent, so it suffices to prove
that $A$ is !!{enumerable} iff $\card{A} \leq \omega$. For
\emph{right-to-left}: if $\card{A} \leq \omega$, then
$\cardle{A}{\omega}$ by
\olref[cardinals][cardsasords]{lem:CardinalsBehaveRight} and
\olref{omegaisacardinal}. For \emph{left-to-right}: suppose $A$ is
!!{enumerable}; then by \olref[sfr][siz][enm-alt]{defn:enumerable} there
are three possible cases:
\begin{enumerate}
	\item if $A = \emptyset$, then $\card{A} = 0 \in \omega$, by
	\olref{naturalsarecardinals} and
	\olref[cardinals][cardsasords]{lem:CardinalsBehaveRight}.
	\item if $\cardeq{n}{A}$, then $\card{A} = n \in \omega$, by
	\olref{naturalsarecardinals} and
	\olref[cardinals][cardsasords]{lem:CardinalsBehaveRight}.
	\item if $\cardeq{\omega}{A}$, then $\card{A} = \omega$, by \olref{omegaisacardinal}.
\end{enumerate}
So in all cases, $\card{A} \leq \omega$. 
\end{proof}
\noindent
Indeed, $\omega$ has a special place. Whilst there are many countable ordinals:

\begin{cor}
$\omega$ is the only !!{enumerable} infinite cardinal.
\end{cor}

\begin{proof}
Let $\cardfont{a}$ be !!a{enumerable} infinite cardinal. Since
$\cardfont{a}$ is infinite, $\omega \leq \cardfont{a}$. Since
$\cardfont{a}$ is !!a{enumerable} cardinal, $\cardfont{a} =
\card{\cardfont{a}} \leq \omega$. So $\cardfont{a} = \omega$ by
Trichotomy. \end{proof}

Of course, there are infinitely many cardinals. So we might ask:
\emph{How many cardinals are there?} The following results show that
we might want to reconsider that question.

\begin{prop}\ollabel{unioncardinalscardinal}
If every member of $X$ is a cardinal, then $\bigcup X$ is a cardinal.
\end{prop}

\begin{proof}
It is easy to check that $\bigcup X$ is an ordinal. Let $\alpha \in
\bigcup X$ be an ordinal; then $\alpha \in \cardfont{b} \in X$ for
some cardinal $\cardfont{b}$. Since $\cardfont{b}$ is a cardinal,
$\cardless{\alpha}{\cardfont{b}}$. Since $\cardfont{b} \subseteq
\bigcup X$, we have $\cardle{\cardfont{b}}{\bigcup X}$, and so
$\cardneq{\alpha}{\bigcup X}$. Generalising, $\bigcup X$ is a
cardinal.
\end{proof} 

\begin{thm}\ollabel{lem:NoLargestCardinal}
There is no largest cardinal.
\end{thm}

\begin{proof}
For any cardinal $\cardfont{a}$, Cantor's Theorem
(\olref[sfr][siz][car]{thm:cantor}) and
\olref[cardinals][cardsasords]{lem:CardinalsExist} entail that
$\cardfont{a} < \card{\Pow{\cardfont{a}}}$.
\end{proof}

\begin{thm}
The set of all cardinals does not exist.
\end{thm}

\begin{proof}
For reductio, suppose $C = \Setabs{\cardfont{a}}{\cardfont{a} \text{
is a cardinal}}$. Now $\bigcup C$ is a cardinal by
\olref{unioncardinalscardinal}, so by \olref{lem:NoLargestCardinal}
there is a cardinal $\cardfont{b} > \bigcup C$. By definition
$\cardfont{b} \in C$, so $\cardfont{b} \subseteq \bigcup{C}$, so that
$\cardfont{b} \leq \bigcup C$, a contradiction.
\end{proof}

You should compare this with both Russell's Paradox and Burali-Forti. 

% ==== OLP file: content/set-theory/cardinals/cp.tex ====


\olfileid{sth}{cardinals}{cp}
\olsection{Cantor's Principle}

Cast your mind back to \olref[ordinals][vn]{sec}. We were discussing
well-ordered sets, and suggested that it would be nice to have objects
which go proxy for well-orders. With this is mind, we introduced
ordinals, and then showed in
\olref[ordinals][ordtype]{ordtypesworklikeyouwant} that these
behave as we would want them to, i.e.:
\[
	\ordtype{A, <} = \ordtype{B, \lessdot} 
	\text{ iff } \tuple{A, <} \isomorphic \tuple{B, \lessdot}.
\]
Cast your mind back even further, to \olref[sfr][siz][equ]{sec}.
There, working na\"ively, we introduced the notion of the ``size'' of
a set. Specifically, we said that two sets are equinumerous,
$\cardeq{A}{B}$, just in case there is !!a{bijection} $f \colon A \to
B$. This is an intrinsically {simpler} notion than that of a
well-ordering: we are only interested in !!{bijection}s, and not (as
with order-isomorphisms) whether the !!{bijection}s ``preserve any
structure''.

This all gives rise to an obvious thought. Just as we introduced
certain objects, \emph{ordinals}, to calibrate well-orders, we can
introduce certain objects, \emph{cardinals}, to calibrate size. That
is the aim of this chapter. 

Before we say what these cardinals will be, we should lay down a
principle which they ought to satisfy. Writing $\card{X}$ for the
cardinality of the set $X$, we would want them to obey:
\[
	\card{A} = \card{B} \text{ iff } \cardeq{A}{B}.
\]
We'll call this \emph{Cantor's} Principle, since Cantor was probably
the first to have it very clearly in mind. (We'll say more about its
relationship to \emph{Hume's} Principle in \olref[hp]{sec}.) So
our aim is to define $\card{X}$, for each $X$, in such a way that it
delivers Cantor's Principle.

% ==== OLP file: content/set-theory/cardinals/hp.tex ====


\olfileid{sth}{cardinals}{hp}
\olsection{Appendix: Hume's Principle}

In \olref[cp]{sec}, we described Cantor's Principle. This was:
\begin{align*}
	\card{A} = \card{B} & \text{ iff } A \approx B.
\intertext{This is very similar to what is now called
\emph{Hume's Principle}, which says:}
	\fregenum{x} {F(x)} = \fregenum{x}{G(x)} & \text{ iff } F \sim G
\end{align*}
where `$F \sim G$' abbreviates that there are exactly as many $F$s as
$G$s, i.e., the $F$s can be put into a bijection with the $G$s, i.e.:
\begin{align*}
	\exists R(&\forall v\forall y(Rvy \lif (Fv \land Gy)) \land {}\\
		&\forall v(Fv \lif \lexists![y][Rvy]) \land {}\\
		&\forall y(Gy \lif \lexists![v][Rvy]))
\end{align*}
But there is a type-difference between Hume's Principle and Cantor's
Principle. In the statement of Cantor's Principle, the variables
``$A$'' and ``$B$'' are first-order terms which stand for \emph{sets}.
In the statement of Hume's Principle, ``$F$'', ``$G$'' and ``$R$'' are
\emph{not} first-order terms; rather, they are in \emph{predicate
position}. (Maybe they stand for \emph{properties}.) So we might gloss
Hume's Principle in English as: the number of $F$s is the number of
$G$s iff the $F$s are bijective with the~$G$s. This is called
\emph{Hume's Principle}, because Hume once wrote this:
\begin{quote}
  When two numbers are so combined as that the one has always an unit
  answering to every unit of the other, we pronounce them equal.
  \citep[Pt.III Bk.1 \S1]{Hume1740}
\end{quote}
And Hume's Principle was brought to contemporary mathematico-logical
prominence by \citet[\S63]{Frege1884}, who quoted this passage from
Hume, before (in effect) sketching (what we have called) Hume's
Principle. 

You should note the structural similarity between Hume's Principle and
Basic Law~V. We formulated this in \olref[story][blv]{sec} as
follows:
\[
	\fregeext{x}{F(x)} = \fregeext{x}{G(x)} \text{iff } \lforall[x][(F(x) \liff G(x))].
\]
And, at this point, some commentary and comparison might help. 

There are two ways to take a principle like Hume's Principle or Basic
Law~V: \emph{predicatively} or \emph{impredicatively} (recall
\olref[story][predicative]{sec}). On the impredicative reading of
Basic Law~V, for each~$F$, the object $\fregeext{x}{F(x)}$ falls
within the domain of quantification that we used in formulating Basic
Law~V itself. Similarly, on the impredicative reading of Hume's
Principle, for each~$F$, the object $\fregenum{x}{F(x)}$ falls within
the domain of quantification that we used in formulating Hume's
Principle. By contrast, on the \emph{predicative} understanding, the
objects $\fregeext{x}{F(x)}$ and~$\fregenum{x}{F(x)}$ would be
entities from some \emph{different} domain. 

Now, if we read Basic Law~V impredicatively, it leads to
inconsistency, via Na\"ive Comprehension (for the details, see
\olref[story][blv]{sec}). Much like Na\"ive Comprehension, it can be
rendered consistent by reading it \emph{predicatively}. But it
probably will not do everything that we wanted it to. 

Hume's Principle, however, \emph{can} consistently be read
impredicatively. And, read thus, it is quite powerful.

To illustrate: consider the predicate ``$x \neq x$'', which obviously
nothing satisfies. Hume's Principle now yields an object $\# x( x\neq
x)$. We might treat this as the number~$0$. Now, on the
\emph{impredicative} understanding---but \emph{only} on the
impredicative understanding---this entity $0$ falls within our
original domain of quantification. So we can sensibly apply Hume's
Principle with the predicate ``$x = 0$'' to obtain an object $\#x (x =
0)$. We might treat this as the number~$1$. Moreover, Hume's Principle
entails that $0 \neq 1$, since there cannot be a bijection from the
non-self-identical objects to the objects identical with $0$ (there
are none of the former, but one of the latter). Now, working
impredicatively again, $1$~falls within our original domain of
quantification. So we can sensibly apply Hume's Principle with the
predicate ``$(x = 0 \lor x = 1)$'' to obtain an object $\#x(x = 0 \lor
x = 1)$. We might treat this as the number~$2$, and we can show that
$0\neq 2$ and $1 \neq 2$ and so on. 

In short, taken impredicatively, Hume's Principle entails that there
are \emph{infinitely many objects}. And this has encouraged
\emph{neo-Fregean logicists} to take Hume's Principle as the
foundation for arithmetic. 

Frege \emph{himself}, though, did not take Hume's Principle as his
foundation for arithmetic. Instead, Frege proved Hume's Principle from
an explicit definition: $\fregenum{x}{F(x)}$ is defined as the extension of
the concept $F \sim \Phi$. In modern terms, we might attempt to render
this as $\fregenum{x}{F(x)} = \Setabs{G}{F \sim G}$; but this will pull us
back into the problems of Na\"ive Comprehension.

% ==== OLP file: content/set-theory/cardinals/milestone.tex ====


\olfileid{sth}{cardinals}{zfc}	
\olsection{$\ZFC$: A Milestone}

With the addition of Well-Ordering, we have reached the final
theoretical milestone. We now have all the axioms required for~$\ZFC$.
In detail:

\begin{defn}
The theory $\ZFC$ has these axioms: Extensionality, Union, Pairs,
Powersets, Infinity, Foundation, Well-Ordering and all instances of
the Separation and Replacement schemes. Otherwise put, $\ZFC$ adds
Well-Ordering to~$\ZF$. 
\end{defn}

$\ZFC$ stands for \emph{Zermelo--Fraenkel} set theory with
\emph{Choice}. Now this might seem slightly odd, since the axiom we
added was called ``Well-Ordering'', not ``Choice''. But, when we later
formulate {Choice}, it will turn out that Well-Ordering is equivalent
(modulo~$\ZF$) to Choice (see \olref[choice][woproblem]{thmwochoice}).
So which to take as our ``basic'' axiom is a matter of indifference.
And the name ``$\ZFC$'' is entirely standard in the literature. 

